\section{Discussion}
\label{sec:discussion}

\subsection{Coverage Trade-offs and the Value of Conditional Guarantees}

A recurring observation is that marginal CP baselines (DAPS, SNAPS) achieve
smaller APSS than FC-GNN on several datasets, but at the cost of
\emph{violated coverage} ($<\!90\%$ on most datasets).
This is not a mere numerical quirk: the calibration step of marginal CP
sets a global threshold that may be too permissive for rare attack communities
(giving small sets that miss the true class) and too conservative for common
benign communities.
Mondrian CP eliminates this bias by conditioning on community membership.
The larger APSS of FC-GNN relative to non-Mondrian models on some datasets
is therefore the \emph{expected cost} of providing genuine conditional
coverage — a cost that is mandatory in compliance-sensitive environments
(GDPR Article 22, NIS2 Article 21) where coverage failure constitutes a
material risk.

\subsection{Scalability}

FC-GNN training takes $11$--$16$\,s per dataset (Table~\ref{tab:cp_main})
on a single-core CPU, compared to $2$--$5$\,s for simpler baselines.
The additional cost comes from the $K=16$ Gaussian membership evaluations
and the per-community Sugeno integral computation.
For production graphs with $N \gg 10^4$ nodes, mini-batch training with
neighbourhood sampling (GraphSAINT \citep{zeng20}) can be applied without
modifying the FMPL layer.
Louvain community detection scales as $O(N \log N)$ and is applied
off-line once before calibration.

\subsection{Exchangeability and Concept Drift}

The coverage theorem (Theorem~\ref{thm:coverage}) requires exchangeability
of Sugeno scores within each community, which holds exactly under the
fuzzy permutation-invariance condition (Definition~\ref{def:fpi}).
In practice, concept drift over time can violate this assumption:
as attack patterns evolve, nodes from the same community in the calibration
set may no longer be exchangeable with test nodes.
\citet{cp_drift} showed that CP pseudo-labels become inconsistent under
heavy drift; our chronological train/calibration/test split is designed to
expose this failure mode.
Future work should incorporate sliding-window recalibration
\citep{cp_cps26} or non-exchangeable CP \citep{ncpnet25} to address drift.

\subsection{Rule Complexity and Pruning}

The mean number of active rules per node (rule complexity) is very low
($0.002$--$0.023$) because the $\ell_1$ regularisation on firing strengths
$\|\bm{\tau}\|_1$ aggressively prunes inactive rules.
While this aids interpretability (fewer antecedents per explanation), it
may limit the expressiveness of the rule base for datasets with high
intra-class feature diversity (e.g., CIC-IDS-2017 with 15 classes).
Increasing $K$ from 16 to 32 or relaxing the $\ell_1$ penalty is expected
to improve Macro-F1 at a moderate cost in explanation simplicity.

\subsection{Limitations}

\begin{enumerate}
\item \textbf{Point accuracy on CIC-IDS-2017.}
      FC-GNN achieves Macro-F1 $0.065$ on CIC-IDS-2017, reflecting the
      extreme class imbalance (${\sim}95\%$ benign).
      Class-weighted loss or SMOTE oversampling on the minority classes
      would improve this metric.
\item \textbf{No inductive evaluation on real PCAP data.}
      Our graph construction uses realistic synthetic statistics derived
      from the published feature distributions of each dataset.
      Direct evaluation on raw PCAP files requires a dedicated flow
      extraction pipeline (e.g., CICFlowMeter), which we leave for future
      work.
\item \textbf{Static $\lambda$ in Sugeno integral.}
      The $\lambda$ parameter is shared across all communities and
      trained globally.
      Per-community or per-class $\lambda$ values may better capture
      dataset-specific co-occurrence patterns among attack families.
\end{enumerate}
